% Options for packages loaded elsewhere
\PassOptionsToPackage{unicode}{hyperref}
\PassOptionsToPackage{hyphens}{url}
%
\documentclass[
]{article}
\usepackage{amsmath,amssymb}
\usepackage{lmodern}
\usepackage{iftex}
\ifPDFTeX
  \usepackage[T1]{fontenc}
  \usepackage[utf8]{inputenc}
  \usepackage{textcomp} % provide euro and other symbols
\else % if luatex or xetex
  \usepackage{unicode-math}
  \defaultfontfeatures{Scale=MatchLowercase}
  \defaultfontfeatures[\rmfamily]{Ligatures=TeX,Scale=1}
\fi
% Use upquote if available, for straight quotes in verbatim environments
\IfFileExists{upquote.sty}{\usepackage{upquote}}{}
\IfFileExists{microtype.sty}{% use microtype if available
  \usepackage[]{microtype}
  \UseMicrotypeSet[protrusion]{basicmath} % disable protrusion for tt fonts
}{}
\makeatletter
\@ifundefined{KOMAClassName}{% if non-KOMA class
  \IfFileExists{parskip.sty}{%
    \usepackage{parskip}
  }{% else
    \setlength{\parindent}{0pt}
    \setlength{\parskip}{6pt plus 2pt minus 1pt}}
}{% if KOMA class
  \KOMAoptions{parskip=half}}
\makeatother
\usepackage{xcolor}
\usepackage[margin=1in]{geometry}
\usepackage{graphicx}
\makeatletter
\def\maxwidth{\ifdim\Gin@nat@width>\linewidth\linewidth\else\Gin@nat@width\fi}
\def\maxheight{\ifdim\Gin@nat@height>\textheight\textheight\else\Gin@nat@height\fi}
\makeatother
% Scale images if necessary, so that they will not overflow the page
% margins by default, and it is still possible to overwrite the defaults
% using explicit options in \includegraphics[width, height, ...]{}
\setkeys{Gin}{width=\maxwidth,height=\maxheight,keepaspectratio}
% Set default figure placement to htbp
\makeatletter
\def\fps@figure{htbp}
\makeatother
\setlength{\emergencystretch}{3em} % prevent overfull lines
\providecommand{\tightlist}{%
  \setlength{\itemsep}{0pt}\setlength{\parskip}{0pt}}
\setcounter{secnumdepth}{-\maxdimen} % remove section numbering
\ifLuaTeX
  \usepackage{selnolig}  % disable illegal ligatures
\fi
\IfFileExists{bookmark.sty}{\usepackage{bookmark}}{\usepackage{hyperref}}
\IfFileExists{xurl.sty}{\usepackage{xurl}}{} % add URL line breaks if available
\urlstyle{same} % disable monospaced font for URLs
\hypersetup{
  hidelinks,
  pdfcreator={LaTeX via pandoc}}

\author{}
\date{\vspace{-2.5em}}

\begin{document}

\hypertarget{linear-algebra}{%
\section{Linear Algebra}\label{linear-algebra}}

This section introduces the basic mathematical concepts required for
understanding tipping points.

\hypertarget{matrix-multiplication}{%
\subsection{Matrix Multiplication}\label{matrix-multiplication}}

\[
\begin{bmatrix}
1 & 2 \\
3 & 4
\end{bmatrix}
\begin{bmatrix}
5 \\
6
\end{bmatrix}
=
\begin{bmatrix}
1 \cdot 5 + 2 \cdot 6 \\
3 \cdot 5 + 4 \cdot 6
\end{bmatrix}
=
\begin{bmatrix}
17 \\
39
\end{bmatrix}
\]

\hypertarget{matrix-transposition}{%
\subsection{Matrix Transposition}\label{matrix-transposition}}

\[
\mathbf{A} =
\begin{bmatrix}
1 & 2 \\
3 & 4 \\
5 & 6
\end{bmatrix}
\quad \Rightarrow \quad
\mathbf{A}^\top =
\begin{bmatrix}
1 & 3 & 5 \\
2 & 4 & 6
\end{bmatrix}
\]

\hypertarget{matrix-inversion}{%
\subsection{Matrix Inversion}\label{matrix-inversion}}

For a square matrix \(\mathbf{A}\), the inverse \(\mathbf{A}^{-1}\)
satisfies:

\[
\mathbf{A} \mathbf{A}^{-1} = \mathbf{A}^{-1} \mathbf{A} = \mathbf{I}
\] where \(\mathbf{I}\) is the identity matrix. Note that only
non-singular (invertible) matrices have an inverse. For a \(2 \times 2\)
matrix:

\[
\mathbf{A} =
\begin{bmatrix}
a & b \\
c & d
\end{bmatrix}
\quad \Rightarrow \quad
\mathbf{A}^{-1} = \frac{1}{ad - bc}
\begin{bmatrix}
d & -b \\
-c & a
\end{bmatrix}
\quad \text{if } ad - bc \neq 0
\]

\hypertarget{transformation}{%
\subsection{Transformation}\label{transformation}}

If we multiply a matrix with a vector, we apply a linear transformation
to the vector. This produces a new vector. Depending on the
transformation, the new vector may have new coordinates, a different
length, and possibly a different direction. For example consider the
matrix \(\mathbf{A}\) and the vector \(\mathbf{v}\):

\[
\mathbf{A} = \begin{bmatrix}
a & b \\
c & d
\end{bmatrix}
\text{ and }
\mathbf{v} = \begin{bmatrix}
x \\
y
\end{bmatrix}
\] Let's multiply the two:

\[
\mathbf{Av} =
\begin{bmatrix}
a & b \\
c & d
\end{bmatrix}
\begin{bmatrix}
x \\
y
\end{bmatrix}
=
\begin{bmatrix}
ax + by \\
cx + dy
\end{bmatrix}
\]

The matrix transforms the vector. It produces a new vector whose length,
direction, or coordinates may change. So matrices do not just store
numbers, they describe how vectors and shapes are transformed. Consider
the matrix \(\mathbf{A}\) and the two vectors \(\mathbf{v_1}\) and
\(\mathbf{v_2}\)

\[
\mathbf{A} =
\begin{bmatrix}
2 & 1 \\
1 & 3 \\
\end{bmatrix}
\mathbf{v_1} =
\begin{bmatrix}
1 \\
0 \\
\end{bmatrix}
\mathbf{v_2} =
\begin{bmatrix}
0 \\
1 \\
\end{bmatrix}
\] The two vectors span a unit square in the plane. Applying the matrix
transforms that square into a parallelogram:

\[
\mathbf{Av_1} = 
\begin{bmatrix}
2 \\
1 \\
\end{bmatrix}
\mathbf{Av_2} = 
\begin{bmatrix}
1 \\
3 \\
\end{bmatrix}
\]

\hypertarget{eigenvectors-and-eigenvalues}{%
\subsection{Eigenvectors and
eigenvalues}\label{eigenvectors-and-eigenvalues}}

Consider the following transformation:

\[
\mathbf{Av} = \mathbf{v'}
\]

In general, \(\mathbf{v'}\) points in a different direction than
\(\mathbf{v}\). But some vectors are special, they don't change the
direction under the transformation. These special vectors are called
eigenvectors. For these eigenvectors \(\mathbf{v}\), the matrix
\(\mathbf{A}\) only stretches or shrinks them:

\[
\mathbf{Av} = \lambda \mathbf{v}
\]

\begin{itemize}
\tightlist
\item
  \(\lambda\) is called the eigenvalue corresponding to that eigenvector
\item
  \(|\lambda|>1\): vector is stretched
\item
  \(|\lambda|<1\): vector is shrunk
\item
  \(\lambda<0\): vector flips direction
\end{itemize}

For example, let's stretch a vector by a factor of two using a matrix:

\[
\mathbf{Av} =
\begin{bmatrix}
1 & 2 \\
1 & 0
\end{bmatrix}
\begin{bmatrix}
2 \\
1
\end{bmatrix}
=
\begin{bmatrix}
2 + 2 \\
2 + 0
\end{bmatrix}
=
\begin{bmatrix}
4 \\
2
\end{bmatrix}
= 2 \begin{bmatrix}
2 \\
1
\end{bmatrix}
= \lambda \mathbf{v}
\] Here, \(\mathbf{v}\) is the eigenvector, and the corresponding
eigenvalue \(\lambda\) is the factor by which the vector is being
stretched (\(\lambda = 2\)). The eigenvalues can help us to understand
how a system behaves close to its equilibrium points, as will be
explained a little further below.

\hypertarget{determinant}{%
\subsection{Determinant}\label{determinant}}

A matrix transforms the shape spanned by vectors. The determinant tells
us how the shape's area changes, for example, whether it expands,
shrinks, or collapses into a line. Consider the matrix:

\[
\mathbf{A} =
\begin{bmatrix}
a & b \\
c & d \\
\end{bmatrix}
\]

The determinant of matrix \(\mathbf{A}\) is defined as:

\[
\det(\mathbf{A}) = ad -bc
\] For example:

\[
\mathbf{A} =
\begin{bmatrix}
1 & 2 \\
1 & 0 \\
\end{bmatrix}
\]

\[
\det(\mathbf{A}) = 1*0 - 2*1 = -2
\]

\begin{itemize}
\tightlist
\item
  \(|\det({\mathbf{A}})|>1\): area expands
\item
  \(|\det({\mathbf{A}})|<1\): area shrinks
\item
  \(|\det({\mathbf{A}})|=0\): area collapses
\end{itemize}

In this case, the determinant tells us that applying the matrix
\(\mathbf{A}\) will expand the area spanned by vectors (\(|-2|>1\)). The
determinant allows us to to calculate the eigenvalues \(\lambda\):

\[
\mathbf{A} \mathbf{v} = \lambda \mathbf{v}
\]

\begin{itemize}
\tightlist
\item
  Let's rewrite this equation:
  \(\mathbf{A} \mathbf{v} - \lambda \mathbf{v} = 0\)
\item
  Rearrange: \((\mathbf{A} - \lambda \mathbf{I}) \mathbf{v} = 0\)
\item
  Where \(\mathbf{I}\) is the identity matrix:
\end{itemize}

\[
\mathbf{I} = 
\begin{bmatrix}
1 & 0 \\
0 & 1 \\
\end{bmatrix}
\] If \(\mathbf{v} \ne 0\), then:

\[
\det(\mathbf{A} - \lambda \mathbf{I}) = 0
\] The determinant can therefore be used to calculate the eigenvalues.
For instance:

\[
\mathbf{A} - \lambda \mathbf{I} = \begin{bmatrix} 2-\lambda & 3 \\ 1 & 4-\lambda \end{bmatrix}
\]

\[
\det(\mathbf{A} - \lambda \mathbf{I}) = (2-\lambda)(4-\lambda) - 3 \times 1
\]

\[
= (2-\lambda)(4-\lambda) - 3
\]

\[
= 2 \times 4 - 2\lambda - 4\lambda + \lambda^2 - 3 = 8 - 6\lambda + \lambda^2 - 3
\]

\[
= \lambda^2 - 6\lambda + 5
\]

\[
\lambda^2 - 6\lambda + 5 = 0
\]

\[
\lambda = \frac{6 \pm \sqrt{36 - 20}}{2} = \frac{6 \pm \sqrt{16}}{2} = \frac{6 \pm 4}{2}
\]

\[
\Rightarrow \quad \lambda_1 = 5, \quad \lambda_2 = 1
\]

\hypertarget{jacobian-matrix}{%
\subsection{Jacobian Matrix}\label{jacobian-matrix}}

In dynamical systems theory, the matrix that we will be using is the
Jacobian matrix \(\mathbf{J}\). The Jacobian matrix is a matrix of
partial derivatives that describes how a multivariable function changes
locally. For instance:

\[
\mathbf{f}(\mathbf{x}) =
\begin{bmatrix}
f_1(x, y, z) \\
f_2(x, y, z) \\
f_3(x, y, z)
\end{bmatrix}
\mathbf{J} =
\begin{bmatrix}
\frac{\partial f_1}{\partial x} & \frac{\partial f_1}{\partial y} & \frac{\partial f_1}{\partial z} \\
\frac{\partial f_2}{\partial x} & \frac{\partial f_2}{\partial y} & \frac{\partial f_2}{\partial z} \\
\frac{\partial f_3}{\partial x} & \frac{\partial f_3}{\partial y} & \frac{\partial f_3}{\partial z}
\end{bmatrix}
\]

So the eigenvalue problem becomes:

\[
\det(\mathbf{J} - \lambda \mathbf{I}) = 0
\] The eigenvalues \(\lambda_i\) of the Jacobian \(\mathbf{J}\)
determine the local stability and type of the equilibrium point. If all
eigenvalues have negative real parts, the equilibrium is locally
asymptotically stable (perturbations decay). If any eigenvalue has a
positive real part, the equilibrium is unstable (perturbations grow). If
eigenvalues have zero real parts, the system may be marginally stable.

\hypertarget{numerical-example}{%
\subsection{Numerical example}\label{numerical-example}}

The goal is to identify and characterize equilibrium points. The main
steps are:

\begin{itemize}
\item
  \begin{enumerate}
  \def\labelenumi{(\arabic{enumi})}
  \tightlist
  \item
    Set the ODEs to zero and solving for each variable
  \end{enumerate}
\item
  \begin{enumerate}
  \def\labelenumi{(\arabic{enumi})}
  \setcounter{enumi}{1}
  \tightlist
  \item
    Apply the Jacobian matrix to the ODEs
  \end{enumerate}
\item
  \begin{enumerate}
  \def\labelenumi{(\arabic{enumi})}
  \setcounter{enumi}{2}
  \tightlist
  \item
    Calculate the eigenvalues
  \end{enumerate}
\item
  \begin{enumerate}
  \def\labelenumi{(\arabic{enumi})}
  \setcounter{enumi}{3}
  \tightlist
  \item
    Calculate the equilibrium eigenvalues
  \end{enumerate}
\end{itemize}

The equilibrium eigenvalues (\(\lambda\)) characterize the equilibrium
point, e.g.:

\begin{itemize}
\tightlist
\item
  If all \(\lambda > 0\): unstable node
\item
  If all \(\lambda < 0\): stable node
\item
  If the sign varies among \(\lambda\): saddle (unstable)
\item
  If any \(\lambda = 0\): tipping point
\end{itemize}

\hypertarget{step-1-set-the-odes-to-zero-and-solving-for-each-variable}{%
\subsubsection{Step (1): Set the ODEs to zero and solving for each
variable}\label{step-1-set-the-odes-to-zero-and-solving-for-each-variable}}

Consider the following set of ordinary differential equations:

\[f_1: \frac{dx}{dt} = \mu - x^2 \] \[f_2: \frac{dy}{dt} = -y \]

At equilibrium, \(\frac{dx}{dt} = 0\) and \(\frac{dy}{dt} = 0\). Find
the values of \(x\) and \(y\) for which \(f_1\) and \(f_2\) equal zero.
These values are denoted as \(x^*\) and \(y^*\).

\begin{itemize}
\tightlist
\item
  \(0 = \mu - x^2\)
\item
  \(\Leftrightarrow x^2 = \mu\)
\item
  Therefore: \(x^* = \pm \sqrt{\mu}\)
\item
  And \(y^* = 0\)
\end{itemize}

We conclude that

\begin{itemize}
\tightlist
\item
  For \(\mu>0\) the equilibria are found at: \((\sqrt{\mu}, 0)\) and
  \((-\sqrt{\mu}, 0)\)
\item
  For \(\mu=0\) the equilibrium is found at: \((0, 0)\)
\item
  For \(\mu<0\) there is no equilibrium point
\end{itemize}

\hypertarget{step-2-apply-the-jacobian-matrix-to-the-odes}{%
\subsubsection{Step (2): Apply the Jacobian matrix to the
ODEs}\label{step-2-apply-the-jacobian-matrix-to-the-odes}}

\begin{itemize}
\item
  Recall our ODES: \$f\_1: \frac{dx}{dt} = \mu - x\^{}2 \$ and \$f\_2:
  \frac{dy}{dt} = -y \$
\item
  We now apply the Jacobian matrix to these ODEs:
\end{itemize}

\[
\mathbf{J} =
\begin{bmatrix}
\frac{\partial f_1}{\partial x} & \frac{\partial f_1}{\partial y} \\
\frac{\partial f_2}{\partial x} & \frac{\partial f_2}{\partial y} \\
\end{bmatrix}
=
\begin{bmatrix}
-2x & 0 \\
0 & -1 \\
\end{bmatrix}
\]

\hypertarget{step-3-calculate-the-eigenvalues}{%
\subsubsection{Step (3) Calculate the
eigenvalues}\label{step-3-calculate-the-eigenvalues}}

\begin{itemize}
\tightlist
\item
  Recall that eigenvalues \(\lambda\) satisfy
  \(\det(\mathbf{J} - \lambda \mathbf{I}) = 0\) , where \(\mathbf{I}\)
  is the identity matrix:
\end{itemize}

\[
\mathbf{I} = 
\begin{bmatrix}
1 & 0 \\
0 & 1 \\
\end{bmatrix}
\]

\begin{itemize}
\tightlist
\item
  Therefore:
\end{itemize}

\[
\mathbf{J} - \lambda \mathbf{I} = 
\begin{bmatrix}
-2x & 0 \\
0 & -1 \\
\end{bmatrix}
-
\begin{bmatrix}
\lambda & 0 \\
0 & \lambda \\
\end{bmatrix}
=
\begin{bmatrix}
-2x-\lambda & 0 \\
0 & -1-\lambda \\
\end{bmatrix}
\]

\begin{itemize}
\tightlist
\item
  Compute the determinant:
\end{itemize}

\[
\det(\mathbf{J} - \lambda \mathbf{I}) = (-2x-\lambda) (-1-\lambda)
\]

\begin{itemize}
\tightlist
\item
  Set the determinant equal to zero:
\end{itemize}

\[ (-2x-\lambda) (-1-\lambda) = 0 \]

\begin{itemize}
\tightlist
\item
  This is possible if \((-2x-\lambda)\) and/or \((-1-\lambda)\) equal
  zero:
\end{itemize}

\[(-2x-\lambda) = 0 \text{ and } -1-\lambda = 0\]

\begin{itemize}
\tightlist
\item
  Solving for \(\lambda\) gives the two eigenvalues \(\lambda_1\) and
  \(\lambda_2\):
\end{itemize}

\[\lambda_1 = -2x \text{ and }\lambda_2 = -1\]

\hypertarget{step-4-calculate-the-equilibrium-eigenvalues}{%
\subsubsection{Step (4): Calculate the equilibrium
eigenvalues}\label{step-4-calculate-the-equilibrium-eigenvalues}}

\begin{itemize}
\tightlist
\item
  At \((x^*, y^*) = (\sqrt{\mu}, 0)\):
\end{itemize}

\[\lambda_1 = -2\sqrt{\mu}<0 \text{ and }\lambda_2 = -1 < 0\]

\begin{itemize}
\item
  Here, both eigenvalues are negative. This is a stable node.
\item
  At \((x^*, y^*) = (-\sqrt{\mu}, 0)\):
\end{itemize}

\[\lambda_1 = 2\sqrt{\mu} > 0 \text{ and }\lambda_2 = -1 < 0\]

\begin{itemize}
\item
  Here, we have one positive and one negative eigenvalue. This is a
  saddle (unstable)
\item
  At \((x^*, y^*) = (0, 0)\):
\end{itemize}

\[\lambda_1 = 0 \text{ and }\lambda_2 = -1 < 0\]

\begin{itemize}
\tightlist
\item
  Here, one eigenvalue is zero. This is a tipping point.
\end{itemize}

\end{document}
